% Options for packages loaded elsewhere
\PassOptionsToPackage{unicode}{hyperref}
\PassOptionsToPackage{hyphens}{url}
\documentclass[
]{article}
\usepackage{xcolor}
\usepackage[margin=1in]{geometry}
\usepackage{amsmath,amssymb}
\setcounter{secnumdepth}{-\maxdimen} % remove section numbering
\usepackage{iftex}
\ifPDFTeX
  \usepackage[T1]{fontenc}
  \usepackage[utf8]{inputenc}
  \usepackage{textcomp} % provide euro and other symbols
\else % if luatex or xetex
  \usepackage{unicode-math} % this also loads fontspec
  \defaultfontfeatures{Scale=MatchLowercase}
  \defaultfontfeatures[\rmfamily]{Ligatures=TeX,Scale=1}
\fi
\usepackage{lmodern}
\ifPDFTeX\else
  % xetex/luatex font selection
\fi
% Use upquote if available, for straight quotes in verbatim environments
\IfFileExists{upquote.sty}{\usepackage{upquote}}{}
\IfFileExists{microtype.sty}{% use microtype if available
  \usepackage[]{microtype}
  \UseMicrotypeSet[protrusion]{basicmath} % disable protrusion for tt fonts
}{}
\makeatletter
\@ifundefined{KOMAClassName}{% if non-KOMA class
  \IfFileExists{parskip.sty}{%
    \usepackage{parskip}
  }{% else
    \setlength{\parindent}{0pt}
    \setlength{\parskip}{6pt plus 2pt minus 1pt}}
}{% if KOMA class
  \KOMAoptions{parskip=half}}
\makeatother
\usepackage{graphicx}
\makeatletter
\newsavebox\pandoc@box
\newcommand*\pandocbounded[1]{% scales image to fit in text height/width
  \sbox\pandoc@box{#1}%
  \Gscale@div\@tempa{\textheight}{\dimexpr\ht\pandoc@box+\dp\pandoc@box\relax}%
  \Gscale@div\@tempb{\linewidth}{\wd\pandoc@box}%
  \ifdim\@tempb\p@<\@tempa\p@\let\@tempa\@tempb\fi% select the smaller of both
  \ifdim\@tempa\p@<\p@\scalebox{\@tempa}{\usebox\pandoc@box}%
  \else\usebox{\pandoc@box}%
  \fi%
}
% Set default figure placement to htbp
\def\fps@figure{htbp}
\makeatother
\setlength{\emergencystretch}{3em} % prevent overfull lines
\providecommand{\tightlist}{%
  \setlength{\itemsep}{0pt}\setlength{\parskip}{0pt}}
\usepackage{bookmark}
\IfFileExists{xurl.sty}{\usepackage{xurl}}{} % add URL line breaks if available
\urlstyle{same}
\hypersetup{
  hidelinks,
  pdfcreator={LaTeX via pandoc}}

\author{}
\date{\vspace{-2.5em}}

\begin{document}

\section{Jeux de Rôle - Résolution de Problèmes
Industriels}\label{jeux-de-ruxf4le---ruxe9solution-de-probluxe8mes-industriels}

Collection de jeux de rôle pédagogiques pour former aux méthodes de
résolution de problèmes en milieu industriel.

\begin{center}\rule{0.5\linewidth}{0.5pt}\end{center}

\subsection{Jeu de Rôle n°1 : Le ``QRQC'' de survie (Réaction à
chaud)}\label{jeu-de-ruxf4le-n1-le-qrqc-de-survie-ruxe9action-uxe0-chaud}

\textbf{Inspiré par :} Le document sur le QRQC (Quick Response Quality
Control)\\
\textbf{Objectif :} Pratiquer le ``Genba'' (aller sur le terrain) et
l'arrêt au premier défaut.

\subsubsection{Le Pitch}\label{le-pitch}

Une ligne d'assemblage de cartes électroniques détecte un composant
monté à l'envers. Le client attend la livraison pour ce soir. La tension
monte.

\subsubsection{Les Personnages (4
joueurs)}\label{les-personnages-4-joueurs}

\begin{enumerate}
\def\labelenumi{\arabic{enumi}.}
\tightlist
\item
  \textbf{L'Opérateur} : Stressé, il veut continuer à produire pour
  tenir les cadences.
\item
  \textbf{Le Chef d'Équipe} : Doit décider s'il arrête la ligne
  (Auto-Qualité).
\item
  \textbf{Le Technicien Qualité} : Armé de sa fiche QQOQCP, il doit
  caractériser le défaut.
\item
  \textbf{Le ``Saboteur'' (Formateur)} : Il donne des informations
  contradictoires si on ne lui pose pas les bonnes questions sur le
  terrain.
\end{enumerate}

\subsubsection{La Règle d'Or (Le
Piège)}\label{la-ruxe8gle-dor-le-piuxe8ge}

Le Chef d'Équipe est tenté de faire un ``tri'' et de continuer.
\textbf{Le piège} : Si l'équipe ne fait pas le ``San Gen Shugi'' (Le
vrai lieu, la vraie pièce, la réalité), ils ne verront pas que le
détrompeur sur la machine est cassé.

\subsubsection{Outils imposés}\label{outils-imposuxe9s}

\begin{itemize}
\tightlist
\item
  QQOQCP
\item
  Principe des 3 réels (Genba, Genbutsu, Genjitsu)
\end{itemize}

\subsubsection{Kit d'indices}\label{kit-dindices}

\begin{itemize}
\tightlist
\item
  📸 Photos de composants montés correctement et à l'envers
\item
  📋 Extrait de la gamme opératoire (version actuelle)
\item
  📋 Extrait de gamme opératoire périmée (avec indication différente)
\item
  🔍 Photo du détrompeur cassé (à ne révéler que si inspection physique)
\item
  📊 Fiche QQOQCP vierge
\end{itemize}

\begin{center}\rule{0.5\linewidth}{0.5pt}\end{center}

\subsection{Jeu de Rôle n°2 : La ``Méthode 8D'' en crise (Résolution
structurée)}\label{jeu-de-ruxf4le-n2-la-muxe9thode-8d-en-crise-ruxe9solution-structuruxe9e}

\textbf{Inspiré par :} Méthode 8D (8 Disciplines)\\
\textbf{Objectif :} Appliquer une démarche structurée face à un problème
récurrent.

\subsubsection{Le Pitch}\label{le-pitch-1}

Un client majeur retourne pour la 3ème fois en 2 mois des lots de pièces
usinées avec des côtes hors tolérances (+0.15mm). Le commercial est
furieux, la production nie toute responsabilité. Le directeur convoque
une équipe d'urgence.

\subsubsection{Les Personnages (5
joueurs)}\label{les-personnages-5-joueurs}

\begin{enumerate}
\def\labelenumi{\arabic{enumi}.}
\tightlist
\item
  \textbf{Le Leader 8D} : Doit guider l'équipe à travers les 8
  disciplines sans sauter d'étapes.
\item
  \textbf{Le Responsable Production} : Défensif, pense que c'est un
  problème de réception client.
\item
  \textbf{Le Méthodes/Process} : Connaît les réglages machines mais est
  débordé.
\item
  \textbf{Le Qualiticien} : A les données de contrôle mais elles sont
  parcellaires.
\item
  \textbf{Le ``Perturbateur'' (Formateur)} : Pousse l'équipe à vouloir
  sauter directement aux actions correctives (D6) sans analyse de cause
  racine (D4-D5).
\end{enumerate}

\subsubsection{La Règle d'Or (Le
Piège)}\label{la-ruxe8gle-dor-le-piuxe8ge-1}

L'équipe aura tendance à vouloir « réparer vite » sans analyser.
\textbf{Le piège} : Les déviations sont intermittentes car l'opérateur
change de tour (il y en a 3) selon la disponibilité, mais un seul tour a
un problème de dilatation thermique non compensé.

\subsubsection{Outils imposés}\label{outils-imposuxe9s-1}

\begin{itemize}
\tightlist
\item
  Diagramme des 8D
\item
  Ishikawa (5M)
\item
  5 Pourquoi
\item
  Plan d'actions avec responsables et délais
\end{itemize}

\subsubsection{Kit d'indices}\label{kit-dindices-1}

\begin{itemize}
\tightlist
\item
  📊 Graphique de contrôle statistique avec points hors limites
  (intermittents)
\item
  📋 Planning d'occupation des 3 tours sur 2 mois
\item
  🔧 Fiches de réglage des 3 tours (avec paramètres légèrement
  différents)
\item
  📸 Photos de pièces OK et NOK avec mesures
\item
  📋 Réclamations clients (dates et n° de lots)
\item
  🌡️ Relevés de température atelier (si demandés)
\item
  📋 Fiche 8D vierge
\end{itemize}

\begin{center}\rule{0.5\linewidth}{0.5pt}\end{center}

\subsection{Jeu de Rôle n°3 : Le ``5 Pourquoi'' du mystère (Recherche de
cause
racine)}\label{jeu-de-ruxf4le-n3-le-5-pourquoi-du-mystuxe8re-recherche-de-cause-racine}

\textbf{Inspiré par :} Méthode des 5 Pourquoi\\
\textbf{Objectif :} Creuser en profondeur pour trouver la vraie cause
racine.

\subsubsection{Le Pitch}\label{le-pitch-2}

Des arrêts intempestifs d'un convoyeur surviennent de manière aléatoire
(3-4 fois par semaine). Le service maintenance a déjà changé 2 fois le
moteur, mais le problème persiste. L'usine perd 2h de production à
chaque arrêt.

\subsubsection{Les Personnages (4
joueurs)}\label{les-personnages-4-joueurs-1}

\begin{enumerate}
\def\labelenumi{\arabic{enumi}.}
\tightlist
\item
  \textbf{L'Animateur 5 Pourquoi} : Doit poser les bonnes questions et
  ne pas accepter les réponses superficielles.
\item
  \textbf{Le Technicien Maintenance} : Frustré d'avoir déjà changé le
  moteur 2 fois.
\item
  \textbf{L'Opérateur Production} : A remarqué que ça arrive souvent
  après le nettoyage du sol.
\item
  \textbf{Le ``Fausseur'' (Formateur)} : Oriente vers des causes
  techniques complexes (roulements, alignement, etc.).
\end{enumerate}

\subsubsection{La Règle d'Or (Le
Piège)}\label{la-ruxe8gle-dor-le-piuxe8ge-2}

L'équipe va naturellement chercher des causes techniques complexes.
\textbf{Le piège} : La vraie cause racine est simple : l'eau de
nettoyage s'infiltre dans le coffret électrique mal fermé, créant des
courts-circuits intermittents.

\subsubsection{Outils imposés}\label{outils-imposuxe9s-2}

\begin{itemize}
\tightlist
\item
  Méthode des 5 Pourquoi (minimum 5 itérations)
\item
  Observation terrain obligatoire
\item
  Ishikawa (5M) en support
\end{itemize}

\subsubsection{Kit d'indices}\label{kit-dindices-2}

\begin{itemize}
\tightlist
\item
  📊 Historique des arrêts (dates et heures)
\item
  📋 Planning de nettoyage de l'atelier
\item
  🔧 Bons de travail maintenance (2 changements moteur)
\item
  📸 Photo du coffret électrique (à ne révéler que si demande
  d'inspection)
\item
  📸 Photo de traces d'eau séchée (indices visuels)
\item
  ⚡ Schéma électrique du convoyeur
\item
  📋 Fiche 5 Pourquoi vierge
\end{itemize}

\begin{center}\rule{0.5\linewidth}{0.5pt}\end{center}

\subsection{Jeu de Rôle n°4 : Le ``Pareto'' sous pression
(Priorisation)}\label{jeu-de-ruxf4le-n4-le-pareto-sous-pression-priorisation}

\textbf{Inspiré par :} Diagramme de Pareto (Loi 80/20)\\
\textbf{Objectif :} Prioriser les actions face à de multiples défauts.

\subsubsection{Le Pitch}\label{le-pitch-3}

Le taux de rebuts de l'atelier peinture a explosé ce mois-ci : 12\% au
lieu des 3\% habituels. La direction exige un plan d'action pour demain.
Il y a des défauts partout : coulures, poussières, manques, mauvaises
teintes, etc.

\subsubsection{Les Personnages (4
joueurs)}\label{les-personnages-4-joueurs-2}

\begin{enumerate}
\def\labelenumi{\arabic{enumi}.}
\tightlist
\item
  \textbf{Le Responsable Qualité} : Doit construire le Pareto et
  identifier les 20\% de causes qui génèrent 80\% des rebuts.
\item
  \textbf{Le Chef d'atelier Peinture} : Veut tout traiter en même temps,
  panique.
\item
  \textbf{Le Technicien Process} : Connaît les paramètres peinture mais
  ne voit pas de changement récent.
\item
  \textbf{Le ``Brouilleur'' (Formateur)} : Insiste sur des défauts
  mineurs mais visibles, détourne l'attention.
\end{enumerate}

\subsubsection{La Règle d'Or (Le
Piège)}\label{la-ruxe8gle-dor-le-piuxe8ge-3}

L'équipe va vouloir s'attaquer à tous les défauts. \textbf{Le piège} :
75\% des rebuts sont dus à un seul type de défaut (poussières), causé
par un filtre de cabine bouché depuis 3 semaines suite à un oubli de
maintenance préventive.

\subsubsection{Outils imposés}\label{outils-imposuxe9s-3}

\begin{itemize}
\tightlist
\item
  Diagramme de Pareto
\item
  Feuille de relevé des défauts
\item
  Plan d'action priorisé
\end{itemize}

\subsubsection{Kit d'indices}\label{kit-dindices-3}

\begin{itemize}
\tightlist
\item
  📊 Liste brute de tous les défauts sur 1 mois (100+ défauts non
  classés)
\item
  📊 Classement des défauts par type (à construire par les joueurs)
\item
  📋 Planning de maintenance préventive (avec un oubli surligné)
\item
  📸 Photo du filtre encrassé
\item
  🔬 Photos de pièces avec différents types de défauts
\item
  📈 Graphique d'évolution du taux de rebut (pic récent)
\item
  📋 Feuille de relevé vierge
\end{itemize}

\begin{center}\rule{0.5\linewidth}{0.5pt}\end{center}

\subsection{Jeu de Rôle n°5 : Le ``Poka-Yoke'' salvateur (Détrompeur
anti-erreur)}\label{jeu-de-ruxf4le-n5-le-poka-yoke-salvateur-duxe9trompeur-anti-erreur}

\textbf{Inspiré par :} Systèmes anti-erreur (Poka-Yoke)\\
\textbf{Objectif :} Concevoir une solution technique empêchant la
réapparition du défaut.

\subsubsection{Le Pitch}\label{le-pitch-4}

Sur une ligne de montage de connecteurs, 15\% des pièces arrivent avec
des joints toriques inversés (couleur rouge au lieu de noir). Le
contrôle final les détecte, mais c'est du gaspillage pur. Les opérateurs
sont formés mais l'erreur persiste.

\subsubsection{Les Personnages (4
joueurs)}\label{les-personnages-4-joueurs-3}

\begin{enumerate}
\def\labelenumi{\arabic{enumi}.}
\tightlist
\item
  \textbf{L'Ingénieur Méthodes} : Doit concevoir un système anti-erreur.
\item
  \textbf{L'Opérateur Expérimenté} : Explique que les joints se
  ressemblent beaucoup, facile de se tromper en cadence.
\item
  \textbf{Le Technicien Maintenance} : Doit valider la faisabilité
  technique du Poka-Yoke.
\item
  \textbf{Le ``Minimaliste'' (Formateur)} : Propose des solutions
  compliquées et coûteuses (capteurs, vision\ldots).
\end{enumerate}

\subsubsection{La Règle d'Or (Le
Piège)}\label{la-ruxe8gle-dor-le-piuxe8ge-4}

L'équipe va chercher des solutions high-tech. \textbf{Le piège} : La
solution optimale est simple et peu coûteuse : des bacs de couleurs
différentes (rouge pour joints rouges, noir pour joints noirs) avec des
couvercles rainurés permettant de prendre qu'un joint à la fois.

\subsubsection{Outils imposés}\label{outils-imposuxe9s-4}

\begin{itemize}
\tightlist
\item
  Grille d'analyse Poka-Yoke
\item
  Matrice coût/efficacité
\item
  Prototype ou schéma de la solution
\end{itemize}

\subsubsection{Kit d'indices}\label{kit-dindices-4}

\begin{itemize}
\tightlist
\item
  📸 Photos des joints (très similaires visuellement)
\item
  📊 Statistiques d'erreurs par poste et par opérateur
\item
  🔧 Photos du poste de travail actuel
\item
  💡 Exemples de Poka-Yoke simples (inspiration)
\item
  📋 Budget disponible pour la solution (500€ max)
\item
  📐 Plans du poste de travail
\item
  📋 Grille d'analyse vierge
\end{itemize}

\begin{center}\rule{0.5\linewidth}{0.5pt}\end{center}

\subsection{Jeu de Rôle n°6 : Le ``Kaizen Blitz'' de l'impossible
(Amélioration continue
express)}\label{jeu-de-ruxf4le-n6-le-kaizen-blitz-de-limpossible-amuxe9lioration-continue-express}

\textbf{Inspiré par :} Kaizen (Amélioration continue)\\
\textbf{Objectif :} Réaliser une amélioration rapide avec les moyens du
bord.

\subsubsection{Le Pitch}\label{le-pitch-5}

Le temps de changement de série (SMED) sur une presse prend 45 minutes.
C'est inacceptable selon le Lean. L'équipe a 2 heures pour réduire ce
temps de 50\% avec un budget de 0€ (uniquement réorganisation et
solutions maison).

\subsubsection{Les Personnages (4-5
joueurs)}\label{les-personnages-4-5-joueurs}

\begin{enumerate}
\def\labelenumi{\arabic{enumi}.}
\tightlist
\item
  \textbf{Le Leader Kaizen} : Anime le Kaizen Blitz, chronomètre en
  main.
\item
  \textbf{Le Régleur} : Celui qui fait habituellement le changement,
  connaît toutes les étapes.
\item
  \textbf{L'Opérateur} : Utilisateur final, a des idées mais n'ose pas
  les dire.
\item
  \textbf{Le Méthodes} : Observe et documente le processus actuel.
\item
  \textbf{Le ``Conservateur'' (Formateur)} : Répète ``on a toujours fait
  comme ça'', résiste au changement.
\end{enumerate}

\subsubsection{La Règle d'Or (Le
Piège)}\label{la-ruxe8gle-dor-le-piuxe8ge-5}

L'équipe va vouloir créer des outils complexes. \textbf{Le piège} : 80\%
du temps est perdu à chercher les outils (non rangés), attendre que la
presse refroidisse (non anticipé), et ajuster manuellement des cotes
(pas de gabarit). Solutions : ombre portée pour outils, préchauffage
pendant production, gabarit de réglage rapide.

\subsubsection{Outils imposés}\label{outils-imposuxe9s-5}

\begin{itemize}
\tightlist
\item
  Chronomètre et découpage opératoire
\item
  Diagramme Spaghetti (flux de déplacement)
\item
  Matrice actions internes/externes
\end{itemize}

\subsubsection{Kit d'indices}\label{kit-dindices-5}

\begin{itemize}
\tightlist
\item
  ⏱️ Relevé détaillé des temps (étape par étape)
\item
  📸 Photos du poste de travail (désorganisé)
\item
  📸 Photos d'exemples de ``shadow board'' (panneaux d'outils)
\item
  📋 Procédure SMED actuelle
\item
  🎥 Vidéo courte du changement de série (option : enregistrement
  simulé)
\item
  📐 Plan de la zone de travail
\item
  📋 Fiche Kaizen vierge
\end{itemize}

\begin{center}\rule{0.5\linewidth}{0.5pt}\end{center}

\subsection{Structure des Kits d'Indices
(Template)}\label{structure-des-kits-dindices-template}

Pour chaque jeu de rôle, préparer :

\subsubsection{📸 Visuels}\label{visuels}

\begin{itemize}
\tightlist
\item
  Photos de pièces bonnes/mauvaises
\item
  Photos du poste de travail
\item
  Photos d'indices cachés (ne révéler que si demande explicite)
\end{itemize}

\subsubsection{📋 Documents}\label{documents}

\begin{itemize}
\tightlist
\item
  Extraits de gammes opératoires
\item
  Procédures (conformes et périmées)
\item
  Fiches de données
\end{itemize}

\subsubsection{📊 Données}\label{donnuxe9es}

\begin{itemize}
\tightlist
\item
  Relevés de mesures
\item
  Historiques
\item
  Statistiques
\item
  Graphiques
\end{itemize}

\subsubsection{🔧 Ressources terrain}\label{ressources-terrain}

\begin{itemize}
\tightlist
\item
  Plans techniques
\item
  Schémas
\item
  Fiches de maintenance
\end{itemize}

\subsubsection{📝 Outils vierges}\label{outils-vierges}

\begin{itemize}
\tightlist
\item
  Fiches QQOQCP
\item
  Matrices 8D
\item
  Diagrammes Ishikawa
\item
  Feuilles de relevé
\end{itemize}

\begin{center}\rule{0.5\linewidth}{0.5pt}\end{center}

\subsection{Conseils d'Animation}\label{conseils-danimation}

\subsubsection{Pour le Formateur}\label{pour-le-formateur}

\begin{enumerate}
\def\labelenumi{\arabic{enumi}.}
\tightlist
\item
  \textbf{Ne jamais donner directement la réponse}
\item
  \textbf{Orienter par des questions} : ``Avez-vous regardé\ldots{} ?'',
  ``Avez-vous vérifié\ldots{} ?''
\item
  \textbf{Laisser l'équipe se tromper} puis débriefer
\item
  \textbf{Chronométrer} pour mettre la pression (réalisme)
\item
  \textbf{Valoriser} les bonnes pratiques (aller sur le terrain, poser
  des questions, utiliser les outils)
\end{enumerate}

\subsubsection{Débriefing (15-20 min)}\label{duxe9briefing-15-20-min}

\begin{itemize}
\tightlist
\item
  Qu'avez-vous appris ?
\item
  Quelles erreurs avez-vous faites ?
\item
  Quels outils ont été efficaces ?
\item
  Comment transposer à votre travail ?
\end{itemize}

\subsubsection{Durée par Jeu}\label{duruxe9e-par-jeu}

\begin{itemize}
\tightlist
\item
  \textbf{Phase jeu} : 30-45 minutes
\item
  \textbf{Débriefing} : 15-20 minutes
\item
  \textbf{Total} : 45-65 minutes
\end{itemize}

\begin{center}\rule{0.5\linewidth}{0.5pt}\end{center}

\subsection{Progression Pédagogique
Recommandée}\label{progression-puxe9dagogique-recommanduxe9e}

\begin{enumerate}
\def\labelenumi{\arabic{enumi}.}
\tightlist
\item
  \textbf{JdR n°1 (QRQC)} : Introduction, réflexe terrain
\item
  \textbf{JdR n°3 (5 Pourquoi)} : Approfondir l'analyse causale
\item
  \textbf{JdR n°4 (Pareto)} : Prioriser efficacement
\item
  \textbf{JdR n°2 (8D)} : Méthodologie complète
\item
  \textbf{JdR n°5 (Poka-Yoke)} : Pérenniser les solutions
\item
  \textbf{JdR n°6 (Kaizen)} : Culture d'amélioration continue
\end{enumerate}

\begin{center}\rule{0.5\linewidth}{0.5pt}\end{center}

\emph{Document créé pour la formation à la résolution de problèmes
industriels}\\
\emph{Pôle Formation UIMM - CVDL}

\end{document}
